\section{模型架构：MoE vs Dense}

\subsection{MoE 的短期优势与长期隐忧}

主持人提出：MoE 和稀疏注意力已成为主流做法，大家对效率的关注是否已超过对效果的关注？

\begin{importantbox}{薛福昭：MoE 长期可能被干掉}
MoE 本质是用 memory（参数量）换 FLOPs——同时意味着 data efficiency 较差。但\textbf{长期来看 FLOPs 永远是最便宜的}：尽管摩尔定律不再完美，芯片提升速度仍比数据增长快得多。因此：
\begin{itemize}
  \item 一切用其他东西换 FLOPs 的方案都是\textbf{短期 solution}
  \item 一切用 FLOPs 换其他东西的方案都是\textbf{长期 solution}
\end{itemize}
模型可能会逐渐回到 Dense，甚至像 Diffusion Model 那样变成 Super-Dense。
\end{importantbox}

\textbf{曹玉（后训练视角）}：完全同意 FLOPs 是唯一确定性变量。但补充一个观点——传统不含合成数据的预训练``可能已经 dead 了''。像 Gemini 这样的模型从预训练阶段就引入大量合成数据，相当于\textbf{用 FLOPs 换 Data}，再用 Data 做 Scaling。如果预训练本身已在做``FLOPs for Data''，那 MoE 作为\textbf{带宽受限硬件体系下的工程解}仍有其价值。

\textbf{刘子伟}：短期和长期存在 Trade-off，个人稍偏 Dense Model。但考虑具体部署场景，MoE 仍有其使用特点。此外，有人在探索下一代计算架构（存储与计算更紧密耦合），这可能在 3--5 年后影响架构选择。

\textbf{谢天宝}：Model Architecture 研究已进入百花齐放状态。MoE 只是工业界的成熟方案，学术界还有更多探索。关键是\textbf{架构设计应跟随 Learning Paradigm 走}——如果未来 Continual Learning 变得重要，架构选择也会随之改变。

\begin{warningbox}{OpenAI 的共识：从 Compute-Bound 到 Data-Bound}
薛福昭指出，OpenAI 在 GPT-4.5 发布时公开表示，整体已从 Compute-Bound 转向 Data-Bound——这是整个 Community 的共识。
\end{warningbox}

\subsection{本章小结}
MoE 是当前硬件约束下的工程最优解，但长期可能回归 Dense。预训练范式已发生变化：合成数据的大量引入使得``传统预训练''的边界模糊化。架构选择最终取决于硬件演进和学习范式。

%% ===== 原生多模态 =====
